% -----------------------------------------------------------------------------
% Description Paragraphs
% -----------------------------------------------------------------------------

\title{Reinforcement Learning for Optimal Setpoint Determination in Energy Utility Systems}
\maketitle

% -----------------------------------------------------------------------------
% Read the text to get a general idea of what is to be written in all three paragraphs
% -----------------------------------------------------------------------------
\par
Energy utility systems are the backbone of modern energy infrastructure. Their complexity ranges from small, localized systems to vast, highly interconnected infrastructure networks. Most of them operate under dynamic conditions where multiple interacting variables influence operational efficiency, cost, reliability, and stability. Conventionally, optimal setpoints for such systems are determined using mathematical optimization. These approaches rely on accurate system models and are typically formulated as constrained optimization problems. While effective for well-defined and moderately sized systems, classical optimization becomes computationally expensive as system complexity increases. Furthermore, time-varying dynamics, uncertainty from the integration of renewable sources, and model-plant mismatches makes it challenging to use purely model-based approaches.

\vspace{0.5pt}
\par
This thesis investigates whether reinforcement learning can find optimal operating setpoints for two case studies of different complexity: a battery energy storage system and a non-linear cookie production factory. For each system, the control task is formulated as a Markov decision process (MDP) with explicit state, action, and reward definitions aligned with operational objectives. The RL agent is trained to produce adaptive setpoint policies that minimize operating cost while satisfying process constraints and preserving stable operation. The learned policies are then benchmarked against established mathematical optimization methods to determine where RL provides practical benefits and where its limitations remain.

\vspace{0.5pt}
\par
The evaluation of the proposed method will be conducted through systematic simulation studies on two energy system case studies of varying complexity. Performance will be assessed using quantitative metrics including total operational cost, constraint violation frequency, and computational time. The results will provide a structured comparison between reinforcement learning and classical optimization and will determine under which conditions RL offers practical advantages for real-time setpoint optimization in complex energy utility systems. 
